\section{The \texttt{NCrystal\_process} McStas Component}
NCrystal\_process component

\subsection*{Identification}
\begin{itemize}
  \item \textbf{Site:} 
  \item \textbf{Author:} NCrystal developers, converted to a Union component by Mads Bertelsen
  \item \textbf{Origin:} NCrystal Developers (European Spallation Source ERIC and DTU Nutech)
  \item \textbf{Date:} 20.08.15
\end{itemize}

\subsection*{Description}
\begin{lstlisting}
This process uses the NCrystal library as a Union scattering process. It only
provides physics; geometry and absorption are handled by the other Union
components. Use Union_make_material to provide the absorption coefficient.
Use Union_make_NCrystal_material for a complete material backed by one
NCrystal configuration.
\end{lstlisting}

\subsection*{Input parameters}
Parameters in \textbf{boldface} are required; the others are optional.

\begin{longtable}{p{0.22\textwidth}p{0.12\textwidth}p{0.46\textwidth}p{0.14\textwidth}}
\toprule
\textbf{Name} & \textbf{Unit} & \textbf{Description} & \textbf{Default} \\
\midrule
\endhead
cfg & str & NCrystal material configuration string (details \textless{}a href="https://github.com/mctools/ncrystal/wiki/Using-NCrystal"\textgreater{}on this page\textless{}/a\textgreater{}). & "" \\
interact\_fraction & 1 & How large a part of the scattering events should use this process 0-1 (sum of all processes in material = 1) & -1 \\
init & string & Name of Union\_init component (typically "init", default) & "init" \\
\bottomrule
\end{longtable}

\subsection*{Links}
\begin{itemize}
  \item Component source code found in file \texttt{NCrystal\_process.comp}.
\end{itemize}
\IfFileExists{NCrystal_process_static.tex}{\input{NCrystal_process_static.tex}}{}
